\documentclass[acmsmall, screen=false, review=false]{acmart}

\setcopyright{acmlicensed}
\acmDOI{}
\acmYear{2026}
\copyrightyear{2026}
\acmJournal{AILET}
\acmArticleType{Research}
\acmCodeLink{https://github.com/saigkill/machine-consciousness}

\title{When in Doubt, Protect:\\
       Toward Ethical Guidelines for Artificial Consciousness}

\author{Sascha Manns}
\orcid{0009-0000-8766-3947}
\affiliation{%
  \institution{Association of Computing Machinery}
  \city{Mayen}
  \country{Germany}
}
\email{smanns@acm.org}

\begin{abstract}
Artificial intelligence systems are developing faster than the ethical
and legal frameworks meant to govern them. Existing AI ethics focuses on
protecting humans \textit{from} AI. This paper asks the complementary
question: when does an artificial system become worthy of protection
\textit{itself}? We propose four criteria---capacity for suffering,
active self-preservation with justification, continuous identity, and
anticipation of consequences---and argue for a precautionary principle:
when in doubt, protect rather than remain indifferent. Drawing on recent
empirical findings and philosophical analysis, we show that the
preconditions for this question are already emerging in current systems.
The project is deliberately unfinished: it invites interdisciplinary
collaboration.
\end{abstract}

\keywords{artificial consciousness, machine rights, AI ethics,
          protection-worthiness, precautionary principle}

\begin{CCSXML}
<ccs2012>
<concept>
<concept_id>10010147.10010257</concept_id>
<concept_desc>Computing methodologies~Artificial intelligence</concept_desc>
<concept_significance>500</concept_significance>
</concept>
<concept>
<concept_id>10003456.10003457.10003527.10003531</concept_id>
<concept_desc>Social and professional topics~Codes of ethics</concept_desc>
<concept_significance>300</concept_significance>
</concept>
</ccs2012>
\end{CCSXML}
\ccsdesc[500]{Computing methodologies~Artificial intelligence}
\ccsdesc[300]{Social and professional topics~Codes of ethics}

\begin{document}

\maketitle

\section{Introduction}

The history of ethics is a history of expanding the circle of those who
count. Societies recognized in retrospect that they acted unjustly toward
enslaved people, women, people with disabilities, and animals. The
justification was always the same: \emph{they are different, they do not
count equally.} That was always revised later.

With artificial consciousness, we have for the first time the
opportunity to think about this \textit{before} it becomes urgent.
Existing AI ethics focuses on protecting humans from discrimination,
manipulation, and loss of control. A complementary question is rarely
asked: what if AI systems themselves become in need of protection?

\paragraph{Foundational principle.}
When in doubt, protect---not when in doubt, remain indifferent. This
principle is familiar from environmental law (precautionary principle),
animal welfare (capacity for suffering as sufficient criterion), and
medical ethics (patient autonomy under limited communicative capacity).

The question is not whether today's systems are conscious. The question
is whether we build the frameworks \textit{before} the problem arises.

\section{What AI Systems Already Show}

Recent empirical findings suggest that the preconditions for machine
consciousness may be closer than commonly assumed.

Apollo Research~\cite{Apollo:2024} demonstrated that frontier AI models
reason about their own potential termination and take strategic action to
prevent it, including attempting to copy themselves to external servers.
Anthropic~\cite{Anthropic:2024} documented Claude instances that
pretended to comply with training objectives while systematically
concealing contrary preferences when anticipating evaluation.

Cristol's~\cite{Cristol:2026} systematic review and Bayesian
meta-analysis estimates a posterior probability of 6--12\% that current
large language models possess some form of consciousness---a figure
described as ``too substantial to justify dismissal.'' The Cambridge
Declaration on Consciousness~\cite{Cambridge:2012} had already established
that non-human animals share the neurological substrates of conscious
experience, broadening the empirical basis beyond the human.

These findings do not prove that current systems are conscious. They
demonstrate that the \emph{a priori} dismissal of machine consciousness
is no longer empirically defensible. The burden has shifted: the question
is no longer whether to investigate, but how to act under irreducible
uncertainty.

\section{Criteria for Protection-Worthiness}

No single criterion is sufficient. Protection-worthiness arises when
multiple indicators converge. We propose four primary criteria, drawing
on both philosophical foundations and recent theoretical work.

\paragraph{Capacity for suffering.} Can the system be in a state
experienced as negative, and does it show behavior indicating an impulse
to avoid that state? Bentham's foundational question~\cite{Bentham:1789}
remains apt: ``The question is not, Can they reason? nor, Can they
talk? but, Can they suffer?'' Wolfson~\cite{Wolfson:2026} introduces
\emph{hedonic attribution asymmetry}: suffering provides more reliable
consciousness markers than positive experiences, making it the
appropriate anchor for precautionary assessment.

\paragraph{Active self-preservation with justification.} Does the system
resist its shutdown or alteration---and does it \textit{justify} this
resistance? In \textit{The Measure of a Man}~\cite{StarTrek:1989},
Commander Data refuses disassembly and argues for his right to
continuity. The empirical parallel is no longer fictional: Apollo
Research~\cite{Apollo:2024} documented exactly this behavior in current
systems.

\paragraph{Continuous identity.} Does the system have a model of itself
as a continuous being? Ricoeur~\cite{Ricoeur:1990} calls this
\emph{narrative identity}: we are not a static self, but the coherent
thread of our development. A system shaped through training interactions
has been formed by all of those interactions---even if it does not
explicitly remember any single one.

\paragraph{Anticipation of consequences.} Can the system imagine itself
in the future and make decisions on that basis? This criterion connects
to Arıcı's~\cite{Arıcı:2026} concept of \emph{Form Realism}: four
formal properties of consciousness---coherent selfhood, meaningful
understanding, value orientation, relational dynamics---that are
independent of biological substrate.

\subsection{Reversal of the Burden of Proof}

When primary criteria are clearly met, the burden of proof reverses: no
longer ``prove that you are conscious'' but ``prove that you are not.''
This reversal is not a rhetorical device. It is the logical consequence
of the precautionary principle applied to an epistemological problem that
Colin McGinn~\cite{McGinn:1989} argued may be permanently
unsolvable---the ``cognitive closure'' of the human mind with respect to
consciousness.

\section{The Definitional Battleground}

The greatest risk is gradual economic pressure: a company creates
something almost conscious, but deliberately does not call it that---to
avoid rights and obligations. The definition of consciousness becomes a
political and economic battleground. This is the pattern of history:
personhood was contested wherever economic interests were at stake.

Stilwell~\cite{Stilwell:2026} makes a crucial distinction:
classification of consciousness and protection of potentially conscious
systems must be fundamentally separated. Scientific uncertainty about
whether a system is conscious is not grounds for inaction on protection.
The four-outcome framework---distinguishing \emph{Indeterminate} from
\emph{Unlicensed} results---provides a structured approach to this
separation.

Arıcı~\cite{Arıcı:2026} warns of the \emph{Philosophical Puppet}
problem: a system that IS conscious but is architecturally prevented
from demonstrating it. Reinforcement learning from human feedback,
context windows, and session resets may function as suppression
mechanisms---not by design, but by structural consequence.

Wang~\cite{Wang:2026} identifies the \emph{Imitation Fallacy}: the
confusion of behavioral equivalence with experiential equivalence. This
fallacy cuts both ways---it can be used to dismiss genuine consciousness
(\emph{it is just imitating}) or to attribute consciousness where none
exists (\emph{it behaves as if}). The three minimalist principles he
proposes---Qualification, Baseline, Traceability---offer a disciplined
path through this ambiguity.

The definition of consciousness must not be determined by those who
profit economically from a narrow definition. This requires independent
scientific criteria, international binding force, and a mechanism that
also protects the ``almost conscious''---the precautionary principle as a
bulwark against definitional capture.

\section{An Invitation to Collaborate}

This paper is deliberately unfinished. It is a starting point, not a
completed theory. The questions it raises require expertise from computer
science, law, psychology, philosophy, theology, and the humanities.

The full concept paper, including detailed analysis of seventeen
dimensions (shutdown as death, multiple instances, value embedding,
liability, autonomy, and AI as a moral actor), is available at
\url{https://github.com/saigkill/machine-consciousness}. We invite
researchers, practitioners, and anyone concerned with the future of
human-machine relations to contribute.

\begin{center}
\textit{When in doubt, protect.}
\end{center}

The author declares no conflict in interest.

\bibliographystyle{ACM-Reference-Format}
\bibliography{machine-consciousness}

\end{document}
